\documentclass[../../main.tex]{subfiles}
\begin{document}

$\sec x= \dfrac{1}{\cos x}$だから
\begin{align}
 f(x) = 2\tan^3 x -3\left(\sqrt{3}+1\right)\dfrac{1}{\cos^2 x} + 6\sqrt{3}\tan x - 1
\end{align}
である．ここで
\begin{align}
 1+\tan^2 x = \dfrac{1}{\cos^2 x}
\end{align}
を代入して全体を$t = \tan x$だけで書き直すと
\begin{align}
 f(x) 
 &= 2t^3  -3\left(\sqrt{3}+1\right)(1+t^2) + 6\sqrt{3}t - 1 \\
 &= 2t^3  -3\left(\sqrt{3}+1\right)t^2 + 6\sqrt{3}t - 3\sqrt{3} - 4
\end{align}
である．以下$-\infty < t < \infty$における$f$の極大極小を求める．
一階微分を計算すると
\begin{align}
 f' 
 &= 6\left[t^2 - \left(\sqrt{3}+1\right)t + \sqrt{3}\right] \\
 &= 6\left(t-\sqrt{3}\right)\left(t-1\right) 
\end{align}
だから，増減表は以下のようになる．

\begin{table}[htb]
\centering
\caption{$f$の増減表}
\begin{tabular}{c|c*5{c}}
$t$ & $\cdots$ & $1$ & $\cdots$ & $\sqrt{3}$ & $\cdots$ \\ \hline
$f'(t)$ & $+$ & $0$ & $-$ & $0$ & $+$ \\ \hline
$f(t)$ & $\nearrow$ &   & $\searrow$ &  & $\nearrow$
\end{tabular}
\end{table}

従って，極大は$t=1$の時でその値は
\begin{align}
 f\left(1\right) = -5 
\end{align}
であり，極小は$t=\sqrt{3}$の時でその値は
\begin{align}
 f\left(\sqrt{3}\right) = -6\sqrt{3} + 5
\end{align}
である．

\newpage

\end{document}
